\subsection{Реализация серверной инфраструктуры}

Серверная часть является фундаментом программной системы и обеспечивает единую точку взаимодействия для обоих клиентских приложений. Все операции~--- создание обращений из мобильного приложения, обработка диспетчером через веб-панель, отслеживание статусов~--- выполняются через один программный интерфейс.

\subsubsection{Реализация базы данных}

Работа с базой данных выполнена с использованием подхода Code-First в Entity Framework Core. Реализованы классы-сущности, конфигурации через Fluent API и контекст базы данных SmartCityDbContext. Конфигурации устанавливают ограничения длины строковых полей, уникальные индексы, правила каскадного удаления и составные ограничения.

Для инкапсуляции логики доступа к данным реализован паттерн <<Репозиторий>>. Специализированный ReportRepository включает методы фильтрации по статусу, получения обращения со связанными данными (категория, медиафайлы, история статусов) и поиска дубликатов по координатам. Репозитории объединены через UnitOfWork~--- класс, обеспечивающий атомарность транзакций и единый контекст базы данных для всех репозиториев в рамках одного запроса.

\subsubsection{Реализация программного интерфейса}

Реализованы контроллеры, обслуживающие запросы от обоих клиентских приложений:
\begin{itemize}
    \item ReportsController~--- CRUD-операции с обращениями, поиск дубликатов по координатам и категории, подтверждение обращения другим пользователем, получение списка обращений с фильтрацией по статусу. Используется мобильным приложением при создании обращений и веб-панелью при диспетчеризации;
    \item CategoriesController~--- получение, создание, редактирование и удаление категорий проблем. Мобильное приложение загружает справочник при создании обращения, веб-панель~--- при управлении справочниками;
    \item ServicesController~--- справочник муниципальных служб, привязка служб к категориям. Используется обоими клиентами для отображения контактных данных ответственных служб;
    \item AuthController~--- регистрация, аутентификация по логину/паролю, аутентификация через Google, получение текущего пользователя, восстановление пароля. Единый механизм аутентификации на основе JWT для обоих клиентов;
    \item EventsController~--- городские события с модерацией.
\end{itemize}

\subsubsection{Ключевые алгоритмы}

Алгоритмы реализованы на стороне сервера и используются обоими клиентскими приложениями через единый программный интерфейс.

\textbf{Алгоритм дедупликации обращений.} При создании нового обращения (как из мобильного приложения, так и через веб-панель) среди активных обращений производится поиск записей с совпадающей категорией в радиусе 50~м от указанной точки. Геопространственный поиск выполняется функцией ST\_DWithin расширения PostGIS. При обнаружении дубликата клиентское приложение предлагает пользователю подтвердить существующее обращение вместо создания нового. На рисунке~\ref{fig:alg-dedup} представлена блок-схема алгоритма.

\begin{figure}[H]
    \centering
    \includegraphics[width=0.5\textwidth]{alg-dedup.png}
    \caption{Алгоритм дедупликации обращений}
    \label{fig:alg-dedup}
\end{figure}


\textbf{Алгоритм обработки смены статуса.} Жизненный цикл обращения включает пять статусов: New, Accepted, InProgress, Resolved, Rejected. При каждом изменении статуса (выполняемом диспетчером через веб-панель) создаётся запись в ReportStatusHistory с указанием даты, времени, идентификатора оператора и комментария. Мобильное приложение отображает историю статусов в карточке обращения и уведомляет пользователя об изменениях. На рисунке~\ref{fig:alg-status} представлена блок-схема.

\begin{figure}[H]
    \centering
    \includegraphics[width=0.5\textwidth]{alg-status.png}
    \caption{Алгоритм обработки смены статуса обращения}
    \label{fig:alg-status}
\end{figure}
